\chapter{Conclusion}
\label{ch:conclusion}

\section{Summary}

This thesis presented an automated system for extracting and formalizing policy rules from academic documents. The key contributions are:

\begin{enumerate}
    \item \textbf{LLM-based Rule Identification}: We demonstrated that large language models, specifically Mistral 7B, can identify policy rules with 99\% accuracy and 0.98 confidence.
    
    \item \textbf{FOL Formalization}: We achieved 100\% success in formalizing policy rules using First-Order Logic with deontic operators, after addressing technical parsing issues.
    
    \item \textbf{SHACL Translation}: We developed a systematic approach to translate FOL formalizations to SHACL shapes for RDF-based validation.
    
    \item \textbf{Logic Expressiveness Analysis}: We provided evidence that SOL and HOL are unnecessary for policy formalization, contradicting assumptions about needing higher expressiveness.
\end{enumerate}

\section{Key Findings}

\subsection{Model Selection}

Smaller, instruction-tuned models (Mistral 7B) outperformed larger models (Llama 3.1 70B) for policy rule classification. This finding has practical implications for deployment cost and efficiency.

\subsection{Logic Sufficiency}

First-Order Logic with deontic operators (O, P, F) is sufficient for institutional policy formalization. The 100\% formalization success rate demonstrates that expressiveness is not the limiting factor.

\subsection{Engineering over Theory}

Practical engineering challenges (JSON parsing, response truncation) proved more significant than theoretical logic expressiveness. The 8.3\% improvement from v1 to v2 came from better parsing, not more expressive logic.

\section{Contributions}

\subsection{Theoretical Contributions}

\begin{itemize}
    \item Evidence that FOL+deontic operators suffice for academic policy formalization
    \item Analysis of why SOL/HOL introduce unnecessary complexity
    \item Framework for systematic FOL-to-SHACL translation
\end{itemize}

\subsection{Practical Contributions}

\begin{itemize}
    \item Complete pipeline: PDF → Rules → FOL → SHACL
    \item Open-source implementation with documented scripts
    \item Evaluation methodology for LLM-based policy extraction
\end{itemize}

\section{Recommendations}

For institutions seeking to automate policy compliance:

\begin{enumerate}
    \item \textbf{Model Selection}: Use Mistral 7B or similar instruction-tuned models
    \item \textbf{Logic Choice}: FOL with deontic operators is adequate
    \item \textbf{Validation Approach}: SHACL provides practical RDF-based validation
    \item \textbf{Engineering Focus}: Invest in robust parsing over complex logic
\end{enumerate}

\section{Conclusion}

This research demonstrates that automated policy rule extraction and formalization is feasible using current LLM and semantic web technologies. The combination of Mistral for identification, FOL for representation, and SHACL for validation provides a complete, practical solution for institutional policy compliance.

The key insight is that theoretical expressiveness (SOL/HOL) is less important than practical engineering (parsing, prompting) for successful policy formalization. Future work should focus on validation against real institutional data and extension to multilingual contexts.
